% Generated draft: requires mathematical and source review.
\documentclass[11pt]{article}
\usepackage[margin=25mm]{geometry}
\usepackage{fontspec}
\setmainfont{Latin Modern Roman}
\usepackage{amsmath,amssymb,mathtools,unicode-math}
\setmathfont{Latin Modern Math}
\usepackage{enumitem,needspace,xurl,hyperref}
\setlength{\parindent}{0pt}
\setlength{\parskip}{5pt}
\setlength{\emergencystretch}{3em}
\newcommand{\N}{\mathbb{N}}
\newcommand{\Z}{\mathbb{Z}}
\newcommand{\Q}{\mathbb{Q}}
\newcommand{\R}{\mathbb{R}}
\newcommand{\C}{\mathbb{C}}
\newcommand{\F}{\mathbb{F}}
\newcommand{\A}{\mathbb{A}}
\newcommand{\PP}{\mathbb{P}}
\newcommand{\E}{\mathbb{E}}
\newcommand{\eps}{\varepsilon}
\newcommand{\dd}{\,\mathrm{d}}
\begin{document}
\section*{UnsolvedMath ID: 2306102}
\textbf{Source problem code:} catalogue entry
\subsection{Definitions and mathematical statement}
For a nonconstant analytic function $g$ on $\mathbb D$, define its valence, under the usual multiplicity-counting convention, by
\[
 V(g)=\sup_{a\in\C}\#\{z\in\mathbb D:g(z)=a\},
\]
where solutions are counted with their zero multiplicities. The value is a positive integer or $\infty$; an infinite supremum need not mean that one particular value has infinitely many preimages. The source does not explicitly fix the multiplicity or infinite-valence conventions.

Let $(v_n)_{n\geq1}$ be a sequence with $v_n\in\{1,2,\ldots\}\cup\{\infty\}$. It is a valence sequence if a single function $f$ analytic on $\mathbb D$ simultaneously satisfies
\[
 V(f^{(n)})=v_n\qquad\text{for every }n\geq1.
\]
The indexing begins with the first derivative; no valence condition on $f$ itself is imposed. Requiring a separate realizing function for each $n$ would be a different, weaker problem.

A. W. Goodman asks for meaningful necessary conditions on $(v_n)$ for this simultaneous realization, and for sufficient conditions ensuring realization. A full characterization would identify the set of all sequences in the displayed value space that admit such an $f$, but the source does not demand one particular form of condition or specify what counts as ``interesting.''

The source permits any analytic $f$, so derivatives can be constant or identically zero. Under the literal preimage-count definition a constant function has infinite valence at its own value; some valence terminology instead excludes constants. This convention must be fixed before constant or polynomial realizers can be included or excluded. No extra growth or nonpolynomial hypothesis is silently imposed.
\subsection{Short English statement}
Which sequences can occur as the exact valences of all successive derivatives of one disk-analytic function? Find useful restrictions and constructive sufficient criteria, retaining infinite entries. The supplied source refers to results of Goodman without identifying them. Its triage says the problem was open in Hayman's 2018 update and no definitive resolution was located; no independent current-status assessment is claimed.
\subsection{Sources}
\begin{itemize}
\item UnsolvedMath permanent ID 2306102, source code catalogue entry, ``Research Problems in Function Theory --- Problem 6.102.'' Registry curation and metadata: Ulam AI and UnsolvedMath Contributors, CC BY 4.0; updated 2026-08-06.
\item Hayman, \emph{Research Problems in Function Theory} (2018 source update), Problem 6.102, attributed to A. W. Goodman: \url{https://arxiv.org/abs/1809.07200}. Supplied source-TeX extraction accessed 2026-07-29. The metadata retains the update for status triage and marks status/release and rights review NEEDS\_REVIEW.
\item Goodman is cited for earlier necessary/sufficient results without a complete bibliographic entry or precise statement. No missing reference is inferred.
\end{itemize}
\end{document}
